\subsection{Recurrent supports on trees}\label{sec:tree-lower}

We first separate elementary contact principles from the special features of
\(\Astar\). A \emph{tagged cycle} means a simple directed cycle of
\(\Phi_V\), and its \emph{support} consists of the spatial vertices and edges
used by that tagged cycle. In particular, the support does not include an
ambient edge merely because both its endpoints occur. A tagged cycle is
\emph{reachable} if some initial tag \((x,0)\), \(x\in V\), eventually enters it.
We call a support \emph{long} if it has at least two edges.

\begin{lemma}[Two-state path support]\label{t:path-support}
For any compulsory-moving controller with at most two states, a primitive tagged cycle whose support is a tree has nonempty path support.  If that path has $a\ge2$ edges, the tagged cycle traverses each edge once in each direction and has period $2a$.  On a one-edge support the tagged period is two or four.
\end{lemma}
\begin{proof}
The path-support conclusion is Corollary~\ref{cor:low-state-tree-supports}.  If the path has at least two edges, an internal vertex $v$ has two incident support edges.  Theorem~\ref{thm:support-capacity} gives
\[
 2\le \sum_{e\ni v}k_e=r_v\le2,
\]
so both incident multiplicities equal one.  Propagation along the connected path gives $k_e=1$ on every edge, and Theorem~\ref{thm:support-capacity} then gives period $2a$.  On a single edge, $r_v=k_e\le2$, so $k_e\in\{1,2\}$ and the period is $2$ or $4$.
\end{proof}

\begin{lemma}[Contact on one tagged cycle]\label{t:same-cycle}
If two trajectories eventually enter the same tagged cycle and its support is a path, they achieve distance-one contact.  Their initial tags and entry times may be arbitrary, and the ambient maze need not be a tree.
\end{lemma}
\begin{proof}
This is Corollary~\ref{cor:same-cycle-path-contact}.
\end{proof}

\begin{lemma}[Contact through an edge support]\label{t:edge-contact}
Suppose one eventual tagged cycle is supported on an edge \(uv\), and the
other eventual tagged cycle visits \(u\) or \(v\). The trajectories achieve
distance-one contact, independently of their entry times and phases.
\end{lemma}
\begin{proof}
After both transients have ended, the first agent is always at \(u\) or
\(v\). At any subsequent visit of the second agent to the shared endpoint,
their distance is at most one. Such visits recur.
\end{proof}

\begin{lemma}[Order inversion]\label{t:order-inversion}
Two synchronous nearest-neighbour walks on a path cannot reverse their
strict spatial order without first achieving distance-one contact.
In particular, contact occurs if their eventual disjoint support intervals
occur in the opposite order from their initial positions.
\end{lemma}
\begin{proof}
In path coordinates the difference of their positions changes by at most
two at each step. As long as its absolute value is at least two, its sign
cannot reverse. This proves the assertion, including the whole transient.
\end{proof}

The support hypothesis in Lemma~\ref{t:same-cycle} and the edge hypothesis
in Lemma~\ref{t:edge-contact} will be established before either lemma is
applied. Neither contact nor equality of spatial positions requires equality
of the internal states.

\subsubsection{A north-boundary invariant}

For a support in a tree, call an edge from a support vertex to a vertex
outside the support a \emph{boundary edge}, directed away from the support.
All directions in the following statements refer to the fixed compass of
Table~\ref{tab:astar}.

\begin{lemma}[Exclusion of a full vertex]\label{t:no-full}
On an induced grid tree, a long tagged cycle of \(\Astar\) contains no
vertex with mask \(NESW\).
\end{lemma}
\begin{proof}
Suppose first that a full vertex \(p\) occurs in state \(1\).
Its next tag is \(s_1\), where \(s\) is its southern neighbour.
If \(s\) is internal to the support, one of its two support exits is north,
and its state-\(1\) exit must continue away from \(p\). The only masks
with two distinct exits including north are \(NE\) and \(NW\).
Accordingly \(s\) has an eastern or western neighbour. That neighbour,
\(s\), \(p\), and the corresponding eastern or western neighbour of \(p\)
form a unit square, contradicting inducedness and the tree hypothesis.
If \(s\) is a support endpoint, returning north from state \(1\) requires
mask \(N\) and returns to \(p_1\). This closes a two-tag cycle, contradicting
long support.

A full vertex occurring only in state \(0\) must be a support endpoint,
by Lemma~\ref{t:path-support}. Its eastern neighbour is internal and is
entered in state \(0\). To continue away from \(p\) in that state and
return west in the other state, this neighbour must have mask \(NW\).
Its northern neighbour and the northern neighbour of \(p\) again complete
a unit square. This excludes the remaining case.
\end{proof}

\begin{lemma}[North boundary]\label{t:north-boundary}
For \(\Astar\) on any finite induced grid tree, every boundary edge of a
long tagged cycle points north.
\end{lemma}
\begin{proof}
An internal support vertex uses both states and selects its two support
edges. The masks with distinct selected directions are
\(NE,ES,NW,EW,SW,NES,NEW,NESW\).
Lemma~\ref{t:no-full} excludes \(NESW\). The mask \(NW\) is also impossible
internally: to pass from its western support neighbour towards its
northern support neighbour, the vertex must be entered in state \(0\).
Every eastward move from a non-full vertex arrives in state \(1\), whereas
the predecessor cannot be full.

The remaining internal possibilities are displayed below. Here \(d/r\)
denotes direction \(d\) and new state \(r\), in the convention
\(\delta(q,M)=(d,r)\). The last column lists only neighbours outside the
two selected support edges.
\begin{center}
\small
\begin{tabular}{@{}llll@{}}
\toprule
Internal mask & From \(0\) & From \(1\) & Possible boundary port\\
\midrule
\(EW,NEW\) & \(\W/0\) & \(\E/1\) & \(N\), for \(NEW\)\\
\(NE\) & \(\N/0\) & \(\E/1\) & none\\
\(SW\) & \(\W/0\) & \(\Sdir/1\) & none\\
\(ES,NES\) & \(\Sdir/1\) & \(\E/1\) & \(N\), for \(NES\)\\
\bottomrule
\end{tabular}
\end{center}
It remains to check support endpoints; no global assumption about the
shape of the path is needed. Let \(c\) be an endpoint and \(v\) its adjacent
internal support vertex.

If \(v\) is east of \(c\), its westward return enters \(c\) in state \(0\).
To avoid immediately repeating that return, the eastward move from \(c\)
must arrive at \(v\) in state \(1\). Only an \(E\) leaf, used in state \(0\),
does so.
If \(v\) is west of \(c\), its eastward return enters \(c\) in state \(1\);
the westward return from \(c\) must arrive at \(v\) in state \(0\).
Exactly \(W\) and \(NW\) in state \(1\) have this behaviour.

If \(v\) is north of \(c\), its southward move arrives at \(c_1\).
The only state-\(1\) northward move is from an \(N\) leaf and arrives in
state \(1\). The internal neighbour cannot be \(SW\), since this would
repeat the same two tags immediately; it is \(ES\) or \(NES\).
Finally, if \(v\) is south of \(c\), then \(v\) has mask \(NE\), whose
northward move enters \(c_0\). The southward return must arrive at \(v\)
in state \(1\).
The candidate masks at \(c\) are \(S,NS,ES,NES,ESW\). Each of the last
three requires an eastern neighbour; together with \(c,v\), and the
eastern support neighbour of \(v\), it forms a unit square. Only
\(S\) and \(NS\) remain.

These four directional equations give the complete endpoint table.
\begin{center}
\small
\begin{tabular}{@{}llll@{}}
\toprule
Endpoint mask & Tag & Direction to support & Boundary port\\
\midrule
\(E\) & \(0\) & \(E\) & none\\
\(N\) & \(1\) & \(N\), towards \(ES/NES\) & none\\
\(S,NS\) & \(0\) & \(S\), towards \(NE\) & \(N\), for \(NS\)\\
\(W,NW\) & \(1\) & \(W\) & \(N\), for \(NW\)\\
\bottomrule
\end{tabular}
\end{center}
Every boundary port in the two tables points north.
\end{proof}

Both preceding lemmas are independent of maze size and of maximum degree.
Inducedness is used when a required fourth side turns a local configuration
into a spatial square.

\subsubsection{Short supports and the six-cell budget}

\begin{lemma}[Edge profiles]\label{t:edge-profiles}
Every edge-supported tagged cycle of \(\Astar\), on any maze, has period
two and belongs to one of the five rows of Table~\ref{t:edge-table}.
On a connected induced tree with more than two cells and maximum degree
at most three, these reduce to the five port profiles in
Table~\ref{t:profile-table}. None has a south-pointing boundary edge.
\end{lemma}
\begin{proof}
A four-tag horizontal cycle would require both states at the western
endpoint to move east. Only an \(E\) leaf does this, and both moves arrive
in state \(1\), which precludes four distinct tags. A four-tag vertical
cycle requires an \(N\) leaf at the lower endpoint. At its upper endpoint
both states must move south, and every such mask resets the state to a
constant. This again precludes four distinct tags.

For period two, solve the two endpoint equations in the semantic table.
For example, a lower state-\(0\) tag that moves north as \(0\) has mask
\(N\) or \(NE\); the upper tag returns south as \(0\) precisely with mask
\(NSW\). This gives the first row. A lower \(0\) that arrives north as
\(1\) requires \(NW\), and its upper \(1\) must return south as \(0\),
again with \(NSW\). A lower \(1\) moving north requires \(N\), and its
upper \(1\) must move south as \(1\), giving the third row. Horizontally,
an eastward arrival as \(0\) requires the full western vertex in state
\(0\), followed by a state-\(0\) westward return; an eastward arrival as
\(1\) followed by return to western state \(0\) requires the western
\(E\) leaf and an eastern \(W\) or \(NW\) tag in state \(1\).
A western tag in state \(1\) would require a westward arrival as \(1\),
which is available only from an eastern \(W\) leaf in state \(0\);
no state-\(1\) eastward move can return to that eastern tag in state \(0\).
These are all possible state pairs and directions.

\begin{table}[!htbp]
\centering
\small
\begin{tabular}{@{}lll@{}}
\toprule
Orientation; endpoint tags & Western/lower mask & Eastern/upper mask\\
\midrule
Vertical; \(0,0\) & \(N,NE\) & \(NSW\)\\
Vertical; \(0,1\) & \(NW\) & \(NSW\)\\
Vertical; \(1,1\) & \(N\) & \(S,NS,SW,ESW,NESW\)\\
Horizontal; \(0,0\) & \(NESW\) & \(EW,NEW,SW\)\\
Horizontal; \(0,1\) & \(E\) & \(W,NW\)\\
\bottomrule
\end{tabular}
\caption{All local edge-supported tagged cycles of \(\Astar\). The two
tags are ordered west--east or south--north.}\label{t:edge-table}
\end{table}

The second row forces western neighbours at both endpoints and therefore
a unit square; it is excluded on an induced tree. The fourth row and
the full-mask option of the third are excluded by the degree bound.
A horizontal \(E\)--\(W\) pair or a vertical \(N\)--\(S\) pair forms an
entire two-cell connected component, excluded by the size assumption.
The remaining possibilities are exactly Table~\ref{t:profile-table}.
\end{proof}

\begin{table}[!htbp]
\centering
\small
\begin{tabularx}{\linewidth}{@{}l l l X c@{}}
\toprule
Profile & Tail/head masks & Tagged cycle & Boundary ports & Cells required\\
\midrule
\(H\) & \(E;NW\) & \(u_0,h_1\) & \(N\) at \(h\) & \(1\)\\
\(D_N\) & \(N;NS\) & \(u_1,h_1\) & \(N\) at \(h\) & \(1\)\\
\(D_W\) & \(N;SW\) & \(u_1,h_1\) & \(W\) at \(h\) & \(1\)\\
\(D_{EW}\) & \(N;ESW\) & \(u_1,h_1\) & \(E,W\) at \(h\) & \(2\)\\
\(Z\) & \(N\) or \(NE;NSW\) & \(u_0,h_0\) &
\(N,W\) at \(h\), and \(E\) at \(u\) if \(u\) has mask \(NE\) & \(2\) or \(3\)\\
\bottomrule
\end{tabularx}
\caption{The five edge profiles on subcubic induced trees with more than
two cells. The tail \(u\) lies west of the head \(h\) in \(H\), and south
of it otherwise. The last column counts required cells outside the
two-cell support, before a connecting path is chosen.}\label{t:profile-table}
\end{table}

\begin{lemma}[Disjointness of recurrent supports]\label{t:disjoint}
Distinct tagged cycles of \(\Astar\) on a subcubic induced grid tree have
disjoint spatial supports.
\end{lemma}
\begin{proof}
An internal support vertex uses both tags and cannot occur on a second
tagged cycle. A shared spatial vertex must therefore be an endpoint of
both supports, in opposite states. The endpoint table in the proof of
Lemma~\ref{t:north-boundary} allows no mask in both endpoint states, so
two long supports cannot intersect.

Comparing that table with Table~\ref{t:profile-table} leaves two possible
mask coincidences for a long and an edge support. An \(NS\) endpoint in
state \(0\) could coincide with a \(D_N\) head in state \(1\), but its
southern neighbour would have to be both an internal \(NE\) vertex and
an \(N\) leaf. An \(N\) endpoint in state \(1\) could coincide with the
lower \(Z\) endpoint in state \(0\), but its northern neighbour would
have to be both \(ES/NES\) and \(NSW\). Neither is possible.
For two edge supports the only opposite-state coincidence is an \(N\)
tail in states \(0\) and \(1\); its upper neighbour would require mask
\(NSW\) for the former and \(NS,SW,\) or \(ESW\) for the latter.
Thus it is impossible as well. The two-cell maze has only its single
edge cycle, so it creates no omitted exception.
\end{proof}

\begin{lemma}[Uniqueness of a long support through six]\label{t:long-unique}
If \(|V|\le6\), \(V\) is a subcubic induced grid tree, and \(\Astar\)
has a long tagged cycle on \(V\), it has no other tagged cycle.
\end{lemma}
\begin{proof}
By Lemma~\ref{t:disjoint}, another support would be disjoint. A long
support has only north boundary edges, whereas any other support has
no south boundary edge, by Lemmas~\ref{t:north-boundary}
and~\ref{t:edge-profiles}. Direct adjacency between the supports is
therefore impossible. Their unique connecting path has an intermediate
vertex.

Two long supports require at least \(3+3+1=7\) vertices. A long support
of at least four vertices paired with an edge support also requires
at least seven. The only remaining budget is a three-vertex long
support, a two-vertex edge support, and one connecting vertex.
An edge profile \(D_{EW}\) or \(Z\) requires another branch off this
minimal union, giving a seventh vertex. For \(H\) or \(D_N\), both
boundary ports are north, so the two-step connector would be \(NS\)
and return to its starting point. For \(D_W\), it must be \(NE\).
The southern support leaf of the \(D_W\) head is then adjacent to the
starting support vertex. It either overlaps the long support or closes
a unit square. Every case is excluded.
\end{proof}

\begin{lemma}[The four short-support configurations]\label{t:four-patterns}
If a subcubic induced grid tree of at most six cells contains two
edge-supported tagged cycles of \(\Astar\), then, up to translation,
it is exactly one of the four six-cell sets in
Table~\ref{t:pattern-coordinates}. Each has exactly two tagged cycles.
\end{lemma}
\begin{proof}
The supports are disjoint by Lemma~\ref{t:disjoint}. Let their unique
connector have \(k\) edges, excluding the two support edges themselves.
The union already has \(k+3\) vertices. A required port off this union
costs a further distinct vertex: identifying that neighbour with the
union would create a spatial cycle. A connector attached to a
\(D_{EW}\) head has one such extra cost. At a \(Z\) head the cost is
one; when it attaches at the eastern port of the lower \(Z\) vertex,
the cost is two. The other profiles have no extra cost after their
unique boundary port has been selected.

The resulting fifteen unordered comparisons are recorded in
Table~\ref{t:pair-table}. Connector words are read from the first named
profile to the second. Exchanging their names only relabels the
connector endpoints and does not transform compass directions.

\begin{table}[!htbp]
\centering
\small
\begin{tabularx}{\linewidth}{@{}l l X@{}}
\toprule
Profile pair & Budget & Outcome\\
\midrule
\(H,H\) & \(k\le3\) & North--north obstruction\\
\(H,D_N\) & \(k\le3\) & North--north obstruction\\
\(D_N,D_N\) & \(k\le3\) & North--north obstruction\\
\(D_W,D_W\) & \(k\le3\) & West--west obstruction\\
\(H,D_W\) & \(k\le3\) & Only \(NEE\); gives \(P_H\)\\
\(D_N,D_W\) & \(k\le3\) & Only \(NEE\); gives \(P_N\)\\
\(H,D_{EW}\) & \(k\le2\) & Square or overlapping supports\\
\(D_N,D_{EW}\) & \(k\le2\) & Square\\
\(D_W,D_{EW}\) & \(k\le2\) & Only \(WW\); gives \(P_E\)\\
\(D_{EW},D_{EW}\) & \(k\le1\) & Adjacent southern leaves close a square\\
\(H,Z\) & \(k\le2\) at head; \(\le1\) below &
Square at head; incompatible lower-port directions\\
\(D_N,Z\) & \(k\le2\) at head; \(\le1\) below &
Square at head; incompatible lower-port directions\\
\(D_W,Z\) & \(k\le2\) at head; \(\le1\) below &
Square at head; direct \(W\) into the lower vertex gives \(P_Z\)\\
\(D_{EW},Z\) & \(k\le1\) at head &
Square; lower attachment exceeds the budget\\
\(Z,Z\) & \(k\le1\) between heads &
No compatible direct ports; other attachments exceed the budget\\
\bottomrule
\end{tabularx}
\caption{All fifteen pairings of the five edge profiles. The proof gives
the direction and square obstruction behind every rejected row. No
compass orientations are identified.}\label{t:pair-table}
\end{table}

Here are the complete route arguments. Two north ports require a
connector beginning \(N\) and ending \(S\). At length at most three,
the only simple possibilities are \(NES,NWS\); their endpoint heads
are horizontally adjacent, so the induced edge closes a square.
Shorter possibilities revisit a vertex. For two west ports the
corresponding words are \(WNE,WSE\), and the vertically adjacent
heads again close a square. A north port joined to a west port
requires final step \(E\). The length-two word \(NE\) closes a
square through the second head's southern leaf. At length three,
\(NNE\) closes the same square one level higher, and \(NEE\) is the
only other simple word. This proves the first six rows.

A \(D_{EW}\) attachment requires a horizontal final step. From a north
port, the budget allows only \(NE\) or \(NW\); the southern support
leaf of the second head is adjacent to the first head, closing a
square or overlapping a support. From a west port, a one-step
connection makes the southern leaves adjacent and closes a square;
among two-step words, \(WE\) repeats the starting vertex and \(WW\)
survives. Two \(D_{EW}\) heads have only the already excluded one-step
geometry. This proves rows seven through ten.

A \(Z\) head accepts outside neighbours north and west, so a connector
entering it ends in \(S\) or \(E\). From a north port with \(k\le2\),
the only simple word is \(NE\), which closes a square through the
lower \(Z\) support vertex. From a west port, \(WE\) repeats the start,
and \(WS\) closes a square through the first head's southern leaf.
Direct westward entry would approach the forbidden eastern side of
the \(Z\) head. At the lower \(Z\) endpoint the outside port is east;
the two extra head neighbours force \(k\le1\), and the one-step
westward connection matches only a \(D_W\) head. A \(D_{EW}\)--\(Z\)
head connection has only one step available: its compatible horizontal
direction again leaves adjacent southern support vertices and hence
a square. Connecting instead to the lower \(Z\) vertex would need
at least seven cells. Finally, a direct \(Z\)--\(Z\) head connector
would have to begin in \(N\) or \(W\) and end in \(S\) or \(E\), which
is impossible in one step; any use of a lower port exceeds the
budget. This proves the remaining five rows.

\begin{table}[!htbp]
\centering
\small
\begin{tabular}{@{}ll@{}}
\toprule
Pattern & Cell set\\
\midrule
\(P_N\) & \(\{(0,-1),(0,0),(0,1),(1,1),(2,1),(2,0)\}\)\\
\(P_E\) & \(\{(-1,0),(0,-1),(0,0),(1,0),(2,0),(2,-1)\}\)\\
\(P_Z\) & \(\{(0,0),(0,1),(0,2),(-1,1),(1,0),(1,-1)\}\)\\
\(P_H\) & \(\{(-1,0),(0,0),(0,1),(1,1),(2,1),(2,0)\}\)\\
\bottomrule
\end{tabular}
\caption{The four possible subcubic six-cell trees with two tagged
cycles. Coordinates retain absolute orientation.}\label{t:pattern-coordinates}
\end{table}

The surviving words and required port cells give exactly the four
displayed sets. Their masks satisfy the stated edge profiles, so
both tagged cycles exist, and no additional branch fits. A third tagged cycle
would have disjoint support and use only the two vertices outside
the first two supports. In \(P_N,P_H\), those two vertices have
masks \(ES,EW\), which do not form any row of Table~\ref{t:edge-table}.
In \(P_E,P_Z\), they are not adjacent. Thus there is no third tagged
cycle, including one supported on connector vertices.
\end{proof}

\subsubsection{Initial-state gates and transient contact}

Two spatially separated tagged cycles do not by themselves constitute a
trap. We must determine which state-\(0\) roots reach them and whether
the approach trajectories have already contacted.

\begin{lemma}[Three closed state-\(0\) gates]\label{t:reset-gates}
In each of \(P_N,P_E,P_Z\), every state-\(0\) root reaches the same
tagged cycle.
\end{lemma}
\begin{proof}
In all three patterns the right cycle is the state-\(1\) vertical edge
at a \(D_W\) head \(r\) and its southern \(N\) leaf \(b\). However,
\(r_0\) moves west as \(0\), and \(b_0\) first enters \(r_0\).
These are escape routes from the right support.

For \(P_N\), put \(\ell=(0,0)\), \(a=(0,-1)\), \(p=(0,1)\),
\(m=(1,1)\), \(r=(2,1)\), and \(b=(2,0)\). The \(NS\) head
\(\ell\) sends both states south into the cycle
\((\ell_1,a_1)\). Every state-\(0\) root reaches it, as witnessed by
\[
 a_0\to\ell_0\to a_1,\qquad
 p_0\to\ell_1,\qquad
 m_0\to p_0,\qquad r_0\to m_0,\qquad b_0\to r_0.
\]
The potentially east-going tag \(p_1\) is never reached by a
state-\(0\) root.

For \(P_E\), put \(\ell=(0,0)\), \(a=(0,-1)\), \(w=(-1,0)\),
\(m=(1,0)\), \(r=(2,0)\), and \(b=(2,-1)\). The \(ESW\) head
\(\ell\) sends both states to \(a_1\), and the left cycle is
\((\ell_1,a_1)\). The remaining roots satisfy
\[
 w_0\to\ell_1,\qquad a_0\to\ell_0,\qquad
 m_0\to\ell_0,\qquad r_0\to m_0,\qquad b_0\to r_0.
\]
Thus every root again enters the left cycle.

For \(P_Z\), write \(z=(0,0)\), \(t=(0,1)\), \(n=(0,2)\),
\(w=(-1,1)\), \(r=(1,0)\), \(b=(1,-1)\).
The left cycle is \((z_0,t_0)\). The leaves \(n,w\) both enter
\(t_1\), and the \(NSW\) rule sends that tag to \(z_0\).
The other roots are on the cycle or satisfy
\[
 r_0\to z_0,\qquad b_0\to r_0.
\]
Consequently the east-going tag \(z_1\) is unreachable from the
distinguished initial state.
\end{proof}

\begin{lemma}[The exceptional path]\label{t:exceptional-path}
All pairs of state-\(0\) roots on \(P_H\) achieve distance-one contact.
\end{lemma}
\begin{proof}
Name its vertices
\[
 w=(-1,0),\quad h=(0,0),\quad p=(0,1),\quad
 m=(1,1),\quad r=(2,1),\quad b=(2,0).
\]
Their induced path order is \(w,h,p,m,r,b\).
The complete twelve-tag map is
\begin{center}
\small
\begin{tabular}{@{}llll@{}}
\toprule
Cell & Mask & From \(0\) & From \(1\)\\
\midrule
\(w\) & \(E\) & \(h_1\) & \(h_1\)\\
\(h\) & \(NW\) & \(p_1\) & \(w_0\)\\
\(p\) & \(ES\) & \(h_1\) & \(m_1\)\\
\(m\) & \(EW\) & \(p_0\) & \(r_1\)\\
\(r\) & \(SW\) & \(m_0\) & \(b_1\)\\
\(b\) & \(N\) & \(r_0\) & \(r_1\)\\
\bottomrule
\end{tabular}
\end{center}
Its two tagged cycles are \(L=(w_0,h_1)\) and \(R=(r_1,b_1)\).
The only state-\(0\) root reaching \(R\) is \(h_0\), along
\(h_0\to p_1\to m_1\to r_1\). Every other root reaches \(L\):
\(w_0\) lies on it and \(p_0,m_0,r_0,b_0\) move successively
left towards it.

Pairs in the same basin contact by Lemma~\ref{t:same-cycle}.
For a cross-basin pair, one root is \(h\). Its only predecessor
in path order is \(w\), which is initially adjacent. Every other
root lies to the right of \(h\), yet its eventual support \(L\)
lies to the left of \(R\). These walks reverse their strict order
and therefore contact by Lemma~\ref{t:order-inversion}.
The argument includes all entry times and the complete transients.
\end{proof}

% The only marked start is h_0. L and R denote cycle supports, not two starts.
\begin{figure}[tbp]
\centering
\begin{tikzpicture}[maze,x=10mm,y=10mm]
 \coordinate (w) at (-1,0); \coordinate (h) at (0,0);
 \coordinate (p) at (0,1); \coordinate (m) at (1,1);
 \coordinate (r) at (2,1); \coordinate (b) at (2,0);
 \draw[quiet] (h)--(p)--(m)--(r);
 \draw[supportA] (w)--(h); \draw[supportB] (r)--(b);
 \node[cell,label=below:{$w$}] at (w) {};
 \node[startA,label=below:{$h_0$}] at (h) {};
 \node[cell,label=above left:{$p$}] at (p) {};
 \node[cell,label=above:{$m$}] at (m) {};
 \node[cell,label=above right:{$r$}] at (r) {};
 \node[cell,label=below:{$b$}] at (b) {};
 \draw[transientA] (.13,.09)--(.13,.86)--(.88,.86);
 \draw[transientA] (1.12,.86)--(1.87,.86);
 \node[font=\footnotesize,anchor=north] at (-.5,-.55) {$L=(w_0,h_1)$};
 \node[font=\footnotesize,anchor=north] at (2,-.55) {$R=(r_1,b_1)$};
 \node[paneltitle] at (3.25,1.05) {One exceptional state-zero root};
 \node[note] at (3.25,.36) {$h_0\to p_1\to m_1\to r_1$ reaches $R$.};
 \node[note] at (3.25,-.33) {All other state-zero roots reach $L$.};
 \mazeCompass{9.1}{.15}
\end{tikzpicture}
\caption{The exceptional path $P_H$, in order $w,h,p,m,r,b$ with $h=(0,0)$. The dashed route is the unique state-zero entry to $R$; initially far pairs in different basins must reverse their spatial order before settling.}
\label{f:exceptional-path}
\end{figure}


\begin{samepage}
\begin{proposition}[Small subcubic trees]\label{t:small-trees}
Two copies of \(\Astar\), both started in state \(0\), achieve
distance-one contact on every induced grid tree of maximum degree
at most three with at most six cells.
\end{proposition}
\begin{proof}
The one- and two-cell cases succeed at time zero. On a larger
finite tree every trajectory eventually enters a tagged cycle,
whose support is a path by Lemma~\ref{t:path-support}.
If a long support exists, its tagged cycle is unique by
Lemma~\ref{t:long-unique}, so Lemma~\ref{t:same-cycle} applies.
If all supports are edges, there is either one tagged cycle,
handled by the same lemma, or the maze is one of the four
configurations of Lemma~\ref{t:four-patterns}.
Lemma~\ref{t:reset-gates} and then Lemma~\ref{t:same-cycle}
handle \(P_N,P_E,P_Z\); Lemma~\ref{t:exceptional-path}
handles \(P_H\).
\end{proof}
\end{samepage}

Trees containing a degree-four vertex are covered by the full-centre
analysis below. The distinction matters: at such a vertex, two
different tagged cycles may share a spatial endpoint in opposite
states, a configuration excluded by Lemma~\ref{t:disjoint} only under
its stated degree bound.
